\section{Conclusion}
\label{sec:conclusion}

We have presented CICADAS, a computational framework that couples the parametric g-formula with
mechanistic PK/PD modeling, disease dynamics, and closed-loop treatment policies to emulate
randomized trials of dynamic anti-seizure strategies from longitudinal observational data. Evaluated
on synthetic ICU cohorts with known ground truth, the framework recovers the generating causal
effect under confounding severe enough to eliminate the apparent treatment benefit, outperforms
standard IPTW and marginal structural model implementations on the same simulated data, and degrades
in characterizable ways under unmeasured confounding, measurement error, and informative censoring.
We have specified the conditions a real dataset must satisfy for the framework to apply---treatment
overlap across the policy range, sufficient measured confounding, dosing variation adequate to
identify PK/PD parameters, a severity readout of adequate fidelity, a characterized censoring
mechanism, and adequately represented subgroups---and we have set out what the framework
structurally cannot do, including evaluating therapies outside current practice.

These results establish methodological feasibility and provide a benchmarked implementation. They do
not establish clinical utility. Because our simulated data and our estimation models share
structural assumptions, and because no real patient data were analyzed, whether CICADAS produces
valid causal estimates in ICU practice remains an open question that only independent external
validation, followed by prospective evaluation of any resulting hypothesis, can answer.
